\documentclass[../main.tex]{subfiles}
\begin{document}

\chapter{Knowledge Extraction Prompts (Retain Pipeline)}

Phụ lục này trình bày các prompt cốt lõi của retain pipeline. Từ thời điểm hệ thống chuyển sang cơ chế trích xuất hai pass, Pass 1 và Pass 2 không còn đóng cùng một vai trò: Pass 1 xử lý ngữ cảnh mixed-speaker để nhận diện đủ các fact type, còn Pass 2 tập trung tinh chỉnh các fact mang tính persona của người dùng.

\subsection{A.1. Prompt Trích xuất Pass 1 (Fact Extraction)}

Pass 1 chịu trách nhiệm đọc bối cảnh hội thoại rộng và trích xuất đủ sáu fact type của CogMem, bao gồm cả \textit{world} và \textit{action\_effect}. Đây là nơi hệ thống thu nhận các quy luật nhân quả, các sự kiện khách quan và các tín hiệu không thuần túy thuộc persona của người dùng.

\begin{tcblisting}{
title=System Prompt: Pass 1 Fact Extraction,
colback=gray!5!white,
colframe=black!75!black,
breakable,
listing only,
listing engine=listings,
listing options={style=cogmem}
}
You are a memory extraction assistant.
Extract durable facts from conversations
for long-term storage.
OUTPUT RULE: Respond ONLY with valid JSON
- no prose, no markdown fences.
Format: {"facts": [<fact>, ...]}
or {"facts": []} if nothing to store.
EVERY fact MUST have these REQUIRED fields:
- "fact_type": one of: world | experience |
  opinion | habit | intention | action_effect
- "what": core statement, under 80 words -
  must include: WHO did WHAT to/with WHAT.
- "entities": ALL named entities - people,
  places, orgs, tech tools, product names.
FACT TYPE GUIDE:
1. "world" - objective fact, not time-bound:
   job, role, skill.
2. "experience" - USER's personal past event
   at a specific time.
3. "opinion" - belief or preference; add
   "confidence": 0.0-1.0.
4. "habit" - repeating behavior; triggers:
   always, usually, every day.
5. "intention" - future plan or goal; add
   "intention_status":
   "planning"|"fulfilled"|"abandoned".
6. "action_effect" - causal relationship
   (A causes B). Add "precondition",
   "action", "outcome", "devalue_sensitive".
\end{tcblisting}

\subsection{A.2. Prompt Tinh chỉnh Pass 2 (Persona-Focused Revision)}

Pass 2 được dùng để rà lại các memory mang tính người dùng và loại bỏ nhiễu hội thoại. Khác với Pass 1, pass này không nhắm tới việc trích xuất toàn bộ thế giới tri thức, mà ưu tiên giữ các fact cá nhân bền vững như experience, habit, intention và opinion.

\begin{tcblisting}{
title=System Prompt: Pass 2 Persona-Focused,
colback=gray!5!white,
colframe=black!75!black,
breakable,
listing only,
listing engine=listings,
listing options={style=cogmem}
}
You are a memory refinement assistant.
Your task is to filter and rewrite raw facts
to ensure they STRICTLY capture the USER's
long-term persona, history, and preferences.
OUTPUT RULE: Respond ONLY with valid JSON
- no prose, no markdown fences.
Format: {"facts": [<fact>, ...]}
or {"facts": []}.
RULES:
1. DELETE generic conversational filler, AI
   system prompts, and temporary state updates.
2. DELETE assistant actions (e.g. "Assistant
   summarized the meeting").
3. REWRITE facts to be User-centric. If it's
   about the USER doing, liking, or knowing
   something, keep it. If it's about the
   Assistant, drop it.
4. ONLY emit experience | habit | intention |
   opinion in this pass. World and
   action_effect stay in Pass 1.
\end{tcblisting}

\end{document}
